\subsection{\texttt{TSUB}}
\noindent\textbf{Bundle encoding.} UOP entry 64, opcode \texttt{0x1067}. Bundle: \texttt{B.OP + B.ARG + B.IOTI}. \texttt{B.OP.Opcode}=\texttt{0x1067}. \texttt{B.IOTI}: selector=\texttt{111 last}, \texttt{SrcTile0}=\texttt{src0}, \texttt{SrcTile1}=\texttt{src1}, \texttt{DstTile}=\texttt{dst}, \texttt{SizeCode}=profile tile size.\par\smallskip
{\small
\paragraph{PTO ISA description.}
Elementwise subtract of two tiles.
\paragraph{Example.}
tile\_a - tile\_b: element at (i,j) = tile\_a[i,j] - tile\_b[i,j].
\par\smallskip
\paragraph{PTO ISA operation.}
For each element \texttt{(i, j)} in the valid region:

\[
\mathrm{dst}_{i,j} = \mathrm{src0}_{i,j} - \mathrm{src1}_{i,j}
\]
}\par\smallskip
\paragraph{Notes.}
FP subtraction follows IEEE 754. Integer subtraction wraps modulo 2\textsuperscript{N}. The VARIANT field selects the rounding mode for FP.
\paragraph{Microcode site table.}
{\scriptsize
\begin{longtable}{@{}R{0.055\linewidth}L{0.23\linewidth}L{0.31\linewidth}L{0.22\linewidth}@{}}
\toprule
Seq & Micro instruction & WO[63:32] fields & WM[31:0] fields \\
\midrule
0 & \texttt{u\_vec.vlds} & \texttt{opc=0x17, x=0, fl=mem, seq=0, kind=b} & \texttt{entry=64, cnt=2, res=vec, var=0} \\
1 & \texttt{u\_vec.vlds} & \texttt{opc=0x17, x=0, fl=mem, seq=1, kind=b} & \texttt{entry=64, cnt=2, res=vec, var=0} \\
2 & \texttt{u\_vec.vsub} & \texttt{opc=0x2B, x=0, fl=alu, seq=2, kind=b} & \texttt{entry=64, cnt=2, res=vec, var=0} \\
3 & \texttt{u\_vec.vsts} & \texttt{opc=0x2A, x=0, fl=mem, seq=3, kind=b} & \texttt{entry=64, cnt=2, res=vec, var=0} \\
\bottomrule
\end{longtable}
}
\paragraph{Microcode body DSL.}
\begin{lstlisting}[style=ptocode]
def uop_tsub(src0, src1, dst):
    dtype = dst.element_type
    valid_rows, valid_cols = dst.valid_shape

    for row in range(0, valid_rows, 1):
        remained = valid_cols
        for col in range(0, valid_cols, pto.get_lanes(dtype)):
            mask, remained = pto.make_mask(dtype, remained)
            lhs = pto.vlds(src0[row, col:])
            rhs = pto.vlds(src1[row, col:])
            subtracted = pto.vsub(lhs, rhs, mask)
            pto.vsts(subtracted, dst[row, col:], mask)
    return
\end{lstlisting}
\Needspace{0.34\textheight}
